\scanpage{164}
\mspage{1 - 37 -}
\margin{n° 413}

\uline{Remarque}.- Si le corps $K$ est $\underset{\sim}{R}$ ou $\underset{\sim}{C}$, et si la variété $X$ est de dimension finie, donc localement compacte, une distribution ponctuelle d'ordre $0$ n'est autre qu'une \uline{mesure} ponctuelle. On vérifiera aisément que la définition du produit de convolution des distributions est dans ce cas en accord avec la définition du produit de convolution des mesures données en Intégration, ch. VIII, § 1. Nous verrons plus tard comment le produit de convolution des mesures et celui des distributions ponctuelles sont deux cas particuliers du produit de convolution des distributions non nécessairement ponctuelles.

12. \uline{Produit de convolution sur une variété de groupe}.

Soit $G$ une variété de groupe de classe $C^r$. Sauf mention expresse du contraire, le produit de convolution de deux distributions ponctuelles sur $G$ sera toujours le produit de convolution associé à la \uline{multiplication} dans $G$. Plus généralement, si $t_i$ $(1\leq i\leq n)$ sont des distributions ponctuelles sur $G$, on désigne par $t_1*\ldots*t_n$ le produit de convolution des $t_i$ relatif à l'application $(x_1,\ldots,x_n)\longrightarrow x_1\ldots x_n$ de $G^n$ dans $G$.

PROPOSITION 20.- \uline{Soient} $t$, $t'$, $t''$ \uline{trois distributions ponctuelles sur} $G$, \uline{d'ordre} $\leq s$, $s'$, $s''$ \uline{respectivement} :

(i) \uline{si} $s+s'+s''\leq r$, \uline{on a} :
\[
t*t'*t''=t*(t'*t'')=(t*t')*t''
\]
(ii) \uline{si} $\varphi$ \uline{est un morphisme de classe} $C^r$ \uline{de} $G$ \uline{dans une autre variété de groupe} $H$ \uline{de classe} $C^r$, \uline{on a pour} $s+s'\leq r$ :
\[
\varphi_*(t*t')=\varphi_*(t)*\varphi_*(t').
\]
Cela résulte aussitôt de la proposition 19 et de son corollaire.
\margin{163}
